\documentclass[]{beamer}
\mode<presentation>

\usepackage[utf8]{inputenc}
\usepackage[portuges]{babel}
\usepackage[T1]{fontenc}
%\usepackage{hyperref}
\usepackage[]{csquotes}

\setbeamertemplate{footline}[page number]
\setbeamercovered{transparent}
\beamertemplatenavigationsymbolsempty

%\def\BLACKBG{true}
\ifdefined\BLACKBG
  \setbeamercolor{block body}{bg=normal text.bg!90!black}
  \setbeamercolor{block body example}{bg=normal text.bg!90!black}
  \setbeamercolor{block title alerted}{use={normal text,alerted text},fg=alerted text.fg!75!normal text.fg,bg=normal text.bg!75!white}
  \setbeamercolor{block title}{bg=black}
  \setbeamercolor{block title example}{use={normal text,example text},fg=example text.fg!75!normal text.fg,bg=normal text.bg!75!white}
  \setbeamercolor{fine separation line}{}
  \setbeamercolor{frametitle}{fg=white}
  \setbeamercolor{item projected}{fg=white}
  \setbeamercolor{normal text}{bg=black,fg=white}
  \setbeamercolor{separation line}{}
  \setbeamercolor{structure}{bg=black,fg=white}
  \setbeamercolor{title}{fg=white}
  \setbeamercolor{titlelike}{fg=white}
\else
  \setbeamercolor{block body}{bg=normal text.bg!90!white}
  \setbeamercolor{block body example}{bg=normal text.bg!90!white}
  \setbeamercolor{block title alerted}{use={normal text,alerted text},fg=alerted text.fg!75!normal text.fg,bg=normal text.bg!75!black}
  \setbeamercolor{block title}{bg=white}
  \setbeamercolor{block title example}{use={normal text,example text},fg=example text.fg!75!normal text.fg,bg=normal text.bg!75!black}
  \setbeamercolor{fine separation line}{}
  \setbeamercolor{frametitle}{fg=black}
  \setbeamercolor{item projected}{fg=black}
  \setbeamercolor{normal text}{bg=white,fg=black}
  \setbeamercolor{separation line}{}
  \setbeamercolor{structure}{bg=white,fg=black}
  \setbeamercolor{title}{fg=black}
  \setbeamercolor{titlelike}{fg=black}
\fi

\newcommand{\invertbgcolor}{
  \ifdefined\BLACKBG
    %\setbeamercolor{normal text}{bg=white,fg=black}
  \fi
}

\newcommand{\backbgcolor}{
  \ifdefined\BLACKBG
    %\setbeamercolor{normal text}{bg=white,fg=black}
  \fi
}


\author[Dmitry]{Dmitry Nix \\ \texttt{@dmitrynix}}

\title{Testes, testes everywhere}
\institute{ /(guru|pug-)pi/i }

\begin{document}
  \begin{frame}
    \titlepage
    % Sobre mim: 
    % - Informática;
    % - Matemática;
    % - GURUPI/PUG-PI;
    % - scoola;
  \end{frame}

  \begin{frame}{su - mario}
    \tableofcontents
    % Dividi a talk em duas:
    %  1. Fundamentação teórica
    %  2. Prática com Ruby, também pode ser chamada de boas práticas
  \end{frame}

  \section{O que é teste automatizado de software?}

  \begin{frame}
    \begin{center}
      \Huge Definições
    \end{center}
  \end{frame}

  \subsection*{Definição de Livro}\label{def1}
  \begin{frame}{\subsecname}
  \blockquote[{\cite[Pressman]{pressman_engenharia_1995}}]{
    […] \emph{Não é incomum}\footnote{Grifo meu} que uma organização de software gaste 40\% do
    esforço de projeto total em teste…
  }
\end{frame}


  \subsection*{Definição da Wikipedia}
  \begin{frame}
  \begin{center}
    \Huge Definição da Wikipedia
  \end{center}
\end{frame}

\begin{frame}{\subsecname}
  \blockquote[{\cite[Wikipedia]{_test_2014}}]{
    In software testing, test automation is the use of special software
    (separate from the software being tested) to control the execution of
    tests and the comparison of actual outcomes with predicted outcomes. Test
    automation can automate some repetitive but necessary tasks in a
    formalized testing process already in place, or add additional testing
    that would be difficult to perform manually.
  }
\end{frame}

\begin{frame}{\subsecname}
  \blockquote[{\cite[Wikipedia]{_test_2014}}]{
    Em teste de software, teste automatizado é um uso especial do software
    (separado do software que será testado) para controlar a execução do testes
    e comparar a saída com a saída que se espera. Teste automatizado pode
    efetuar algumas tarefas repetitivas, [\ldots]
  }
\end{frame}

\begin{frame}{}
  \begin{center}
    \includegraphics[width=5cm]{images/like-a-sir}
  \end{center}
\end{frame}


  \subsection*{Minha definição}
  \begin{frame}{\subsecname}
  \begin{block}{Teste}
    É algo que testo =/
  \end{block}
  \pause

  \begin{block}{Automatizado}
    É algo que automatizo :P
  \end{block}
  \pause

  \begin{block}{Software}
    Conjunto de instruções.
  \end{block}
  \pause
\end{frame}

\begin{frame}{\subsecname}
  \begin{center}
    \includegraphics[width=5cm]{images/pokerface}
  \end{center}
\end{frame}


  \backbgcolor
  \section{Exemplos e casos de uso com Ruby}

  \begin{frame}{fibonacci}
    Sequência de Fibonacci :)
  \end{frame}

  \bibliographystyle{plain}
  \bibliography{Items}

  \begin{frame}{}
    \begin{center}
      \includegraphics[width=5cm]{images/isso-e-tudo}
    \end{center}
    % http://lokotopia.com.br/wp-content/uploads/2013/06/013-stefan-asafti-marcas-cartoon-isso-e-tudo.jpg
  \end{frame}
\end{document}
